\section{The Page That Fits in One Service}
\label{sec:ch13-the-page-that-fits}
\index{pagination!across a subgraph boundary|(}%

\code{Product.reviews} is the only field in Mosaic that pages something one
service owns off an entity another service owns. It has been that shape since
chapter~\ref{ch:schema-design-that-survives-change} and it survived the
extraction without anybody thinking about it, so I want to start by saying what
it costs, which is nothing.

\begin{minted}[fontsize=\footnotesize,breaklines]{graphql}
{ browseProducts(first: 2) { nodes { title
    reviews(first: 2) { totalCount pageInfo { endCursor hasNextPage } nodes { rating } } } } }
\end{minted}

Each product comes back with a page of its own, its own \code{totalCount} and
its own \code{endCursor}. The first two products in the catalogue report six
reviews and three, and their cursors are different strings.

\index{pagination!why the reviews connection works}%
The reason it works is that nothing about the page crosses a boundary. The
router locates \code{Product} in Catalog, hands Reviews a representation, and
Reviews runs one keyset query per branch of its DataLoader, which is
chapter~\ref{ch:schema-design-that-survives-change}'s arrangement untouched. The
items, the cursors, the count and the ordering are four answers from one service
about one table.
The boundary is above the connection, not through it. That distinction is the
whole of section~\ref{sec:ch13-the-page-you-cannot-ask-for}, and it is easier to
see now, while it is still working.

\subsection*{What a cursor actually is}

\index{cursor!what it encodes}%
It is worth decoding one, because the rest of this section depends on what is
inside. Take a cursor off a page of products and base64-decode it:

\begin{minted}[fontsize=\scriptsize,breaklines]{text}
e31CZWVjaCBDdXR0aW5nIEJvYXJkOmEwMDAwMDAwLTAwMDAtNDAwMC04MDAwLTAwMDAwMDAwMDAxNg==
\end{minted}

\begin{minted}[fontsize=\footnotesize,breaklines]{text}
{}Beech Cutting Board:a0000000-0000-4000-8000-000000000016
\end{minted}

\index{pagination!keyset rather than offset}%
That is the sort key, a colon, and the tiebreaker.
\code{CatalogService.DefaultOrder} sorts by \code{Title} and appends \code{Id},
and both are in there. Which tells you the important thing about HotChocolate
16's connections: they are \textbf{keyset} paged, not offset paged. The cursor
does not say \enquote{you have had twenty rows}. It says \enquote{you stopped
just after this one}, and the next page is a \code{WHERE} clause rather than an
\code{OFFSET}.

The same holds for reviews, whose order is \code{OrderBy(r => r.CreatedAt)}:

\begin{minted}[fontsize=\footnotesize,breaklines]{text}
{}202505311348000000000+0000:e0000000-0000-4000-8000-000000000087
\end{minted}

Ascending, so a page of reviews is oldest first. The doc comment on that
resolver said \enquote{newest page first} until this chapter, which is the kind
of thing that stays wrong for five chapters because nobody reads a comment next
to a field that works.

\subsection*{A row arriving ahead of an open cursor}

\index{cursor!stability under an insert}%
Chapter~\ref{ch:schema-design-that-survives-change} left a question here: what
happens to a client holding a cursor when a row is inserted ahead of it. Now
there is an answer, and it is the good one.

Take page one of a product's reviews, then insert a review into the Reviews
database dated one second before the row the cursor points at, then ask for page
two:

\begin{minted}[fontsize=\footnotesize,breaklines]{text}
page 1                          page 2, after the insert
2025-05-31T13:48:00Z  rating 5  2025-10-09T10:12:00Z  rating 5
2025-08-20T17:35:00Z  rating 4  2026-01-05T19:26:00Z  rating 4
totalCount 6                    totalCount 7
\end{minted}

Nothing repeated and nothing skipped that the client had not already gone past.
Ask the same database the same question by position and it is a different story:

\begin{minted}[fontsize=\footnotesize,breaklines]{sql}
select created_at, rating from reviews where product_id = ...
  order by created_at offset 2 limit 2;
\end{minted}

\begin{minted}[fontsize=\footnotesize,breaklines]{text}
2025-08-20 17:35:00+00|4
2025-10-09 10:12:00+00|5
\end{minted}

The first of those is the row page one already showed. That is what offset
paging does to a list that is being written to, and it is the reason
chapter~\ref{ch:schema-design-that-survives-change} took a breaking change to
introduce connections in the first place.

The keyset version is not free of consequences either. The inserted review is
invisible to that client until it starts again from page one, because the cursor
has already gone past where the row landed. A client is being shown a stable
view of a list that has moved, which is a different thing from a correct view
and is usually what you want.

\subsection*{The defect that was hiding behind the router}

\index{pagination!the cursor that could not tell two rows apart|(}%
While measuring all that, I found something wrong with the other connection.
Everything in this subsection reproduces at tag \code{ch12} and not at
\code{ch13}, because the last part of it is the repair.

Check out \code{ch12}, ask Catalog for a page of products, and select nothing
but the title:

\begin{minted}[fontsize=\footnotesize,breaklines]{graphql}
{ browseProducts(first: 2) { edges { cursor node { title } } } }
\end{minted}

The cursors decode to this:

\begin{minted}[fontsize=\footnotesize,breaklines]{text}
{}Beech Cutting Board:00000000-0000-0000-0000-000000000000
{}Brant Ottoman:00000000-0000-0000-0000-000000000000
\end{minted}

Every tiebreaker is an empty \code{Guid}. Add \code{id} to the selection set and
they are real identifiers again. And a cursor that cannot tell two rows apart
does exactly what you would expect on the next page:

\begin{minted}[fontsize=\footnotesize,breaklines]{text}
page 1: Beech Cutting Board | Brant Ottoman
page 2: Brant Ottoman | Brant Two-Seat Sofa
\end{minted}

The client is handed \enquote{Brant Ottoman} twice. This is not a rare
interleaving or a race. It is every second page of every query that does not
happen to ask for the identifier, and it has been in the repository since
chapter~\ref{ch:data-without-the-n-plus-1} created the field.

\index{projection!and the sort key}%
The cause is two mechanisms that each work and do not know about each other. The
sort is applied to the query: \code{DefaultOrder} appends \code{Id} and the
\code{ORDER BY} really does carry it. The projection is built somewhere else
entirely, out of the client's selection set, by
\code{QueryContextParameterExpressionBuilder}:

\begin{minted}[fontsize=\footnotesize,breaklines]{csharp}
return new QueryContext<T>(
    selection.AsSelector<T>(context.IncludeFlags),
    filterContext?.AsPredicate<T>(),
    sortContext?.AsSortDefinition<T>());
\end{minted}

Three things built from three sources, and nothing adds the sort keys to the
selector. So EF Core loads the two columns the client asked about, materialises
a \code{Product} whose \code{Id} is default, and the cursor is serialised from
that object rather than from the row.

\subsection*{Why eight chapters of federated queries never found it}

Here is the part that belongs in this book rather than in a HotChocolate issue.
Through the router, the same page is correct as soon as the query asks for
anything at all from another subgraph:

\begin{minted}[fontsize=\footnotesize,breaklines]{text}
selection                            page 2
title                                Brant Ottoman | Brant Two-Seat Sofa
title price { amount }               Brant Two-Seat Sofa | Canvas Storage Bin
title availableQuantity              Brant Two-Seat Sofa | Canvas Storage Bin
title reviews(first: 1){ totalCount} Brant Two-Seat Sofa | Canvas Storage Bin
\end{minted}

\index{query plan!adds the key to a fetch}%
The planner needs \code{Product.id} to build representations, so it adds it to
the Catalog fetch, which puts the key back in the projection, which fixes the
cursor. The only query that reproduces the defect is one that stays inside a
single subgraph, and this book has spent five chapters demonstrating queries
that do not.

I want to be careful about what that is evidence of. It is not evidence that
federation is safer. It is evidence that a query plan changes the selection set
a subgraph sees, so a subgraph tested through a router has been tested against a
different query than the one a client wrote. That is worth knowing about long
after this particular defect stops being interesting.

\subsection*{One line, and a constructor}

The fix is a supported API and it is one line:

\begin{minted}[fontsize=\footnotesize,breaklines]{csharp}
return await db.Products
    .AsNoTracking()
    .With((query ?? QueryContext<Product>.Empty).Include(p => p.Id), DefaultOrder)
    .ToPageAsync(pagingArguments, cancellationToken);
\end{minted}

\code{QueryContext.Include} adds a property to the selector whatever the client
selected. It also does not work, first time, on any of Mosaic's domain records,
and the reason is worth a paragraph because the failure says nothing.
\code{Include} builds its selector by calling \code{ExpressionHelpers.Rewrite},
which ends with

\begin{minted}[fontsize=\footnotesize,breaklines]{csharp}
var newInitExpression = Expression.MemberInit(Expression.New(typeof(TRoot)), bindings);
\end{minted}

and \code{Expression.New(Type)} wants a parameterless constructor.
\code{Product} is a positional record and has none, so the expression throws
while the resolver is being executed and the client gets
\code{Unexpected Execution Error} with nothing in the log. \code{Product} gained
a parameterless constructor whose only caller is that expression, and whose
doc comment says so.

\index{schema!a fix no schema can see}%
Afterwards the query is still projected, which was the point of doing it this
way rather than dropping the selector:

\begin{minted}[fontsize=\footnotesize,breaklines]{sql}
SELECT p.title AS "Title", p.id AS "Id"
FROM products AS p
ORDER BY p.title, p.id
LIMIT @p
\end{minted}

Two columns out of five, and the tiebreaker among them. And the published schema
is byte-identical before and after. Re-export \code{catalog.graphql} and
\code{git diff} says nothing changed, so no schema-comparison tool, no
composition check and no breaking-change gate could have seen either the defect
or the repair. Chapter~\ref{ch:schema-design-that-survives-change} made the same
point about \code{[ID]} doing nothing for four chapters while the SDL said
\code{id: ID!} throughout. The schema is a contract about shapes, and both of
these were faults in what the shape was filled with.

\medskip
So a page inside one service can be got right. The next question is the one a
client asks about ten minutes after seeing this graph, and the answer to that one
is no.

\index{pagination!the cursor that could not tell two rows apart|)}%
\index{pagination!across a subgraph boundary|)}%
